% Source: https://artofproblemsolving.com/community/c3h2325529_geometry
% Stable identity: aops-topic-2325529-q01
\textbf{Shape of the triangle.} Let $AC=x$, so $AB=2x$. By the law of cosines, $BC^2=4x^2+x^2-2\cdot2x\cdot x\cdot\dfrac12=3x^2$, so $BC=x\sqrt3$ and $AC^2+BC^2=4x^2=AB^2$: the angle at $C$ is right (a $30$-$60$-$90$ triangle).

\textbf{Coordinates.} Put $C=(0,0)$, $A=(x,0)$, $B=\left(0,x\sqrt3\right)$ and $P=(p,q)$. Then
\[p^2+q^2=4,\qquad(p-x)^2+q^2=3,\qquad p^2+\left(q-x\sqrt3\right)^2=25.\]
Subtracting the first equation from the other two gives $x^2-2px=-1$ and $3x^2-2\sqrt3\,qx=21$, so
\[p=\dfrac{x^2+1}{2x},\qquad q=\dfrac{3x^2-21}{2\sqrt3\,x}.\]
Substituting into $p^2+q^2=4$ and multiplying by $12x^2$:
\[3(x^2+1)^2+(3x^2-21)^2=48x^2\iff12x^4-168x^2+444=0\iff x^4-14x^2+37=0,\]
so $x^2=7\pm2\sqrt3$.

\begin{asy}
usepackage("amsmath");
size(6.5cm);
defaultpen(fontsize(10pt));
real x=sqrt(7+2*sqrt(3));
pair C=(0,0), A=(x,0), B=(0,x*sqrt(3));
real p=(x^2+1)/(2x), q=(3x^2-21)/(2*sqrt(3)*x); pair P=(p,q);
draw(A--B--C--cycle,linewidth(0.9)); draw(P--A,blue); draw(P--B,red); draw(P--C,dashed);
draw((0.25,0)--(0.25,0.25)--(0,0.25));
dot(A); dot(B); dot(C); dot(P);
label("$A$",A,SE); label("$B$",B,N); label("$C$",C,SW); label("$P$",P,dir(75));
label("$\sqrt3$",P+0.4*(A-P),SW,blue+fontsize(8pt)); label("$5$",(P+B)/2,E,red); label("$2$",(P+C)/2,NW);
label("$x$",(C+A)/2,S); label("$x\sqrt3$",(C+B)/2,W);
\end{asy}

\textbf{Choosing the root.} If $x^2=7-2\sqrt3\approx3.54$, then $AB=2x\approx3.76$ and $BC\approx3.26$; but every point of the triangle is within $\max(BA,BC)$ of $B$, while $PB=5$. So $x^2=7+2\sqrt3$. (Then $p\approx1.77$, $q\approx0.93$ and $\dfrac px+\dfrac{q}{x\sqrt3}\approx0.71<1$, so $P$ is indeed inside.)

\textbf{Area.} The legs are $x$ and $x\sqrt3$, so
\[[ABC]=\dfrac{\sqrt3}{2}x^2=\dfrac{\sqrt3}{2}\left(7+2\sqrt3\right)=\boxed{\dfrac{6+7\sqrt3}{2}}\approx9.06.\]
